
Automated research pipelines fail when hypotheses violate feasibility constraints during phase transitions. When hypotheses pass early validation but reach Phase 4 implementation to discover they require unavailable datasets or human raters, all prior work is lost. Existing workflow systems rely on schema validation (e.g., Pydantic typed fields) to check structured outputs at phase boundaries. Schema validation verifies structure---types are correct, required fields present---but cannot express semantic or compositional constraints.

Consider a research pipeline with feasibility constraints: ''no synthetic data generation,'' ''no human evaluation,'' ''standard datasets only,'' ''existing benchmarks only.'' Schema validation checks that dataset\_name is a string but cannot verify it lacks keywords like ''synthetic'' or ''generated.'' It validates that both dataset\_type and dataset\_name fields exist but cannot enforce the compositional rule ''if dataset\_type equals 'standard' then dataset\_name must be in the approved standard datasets list.''

This expressiveness gap has consequences. Workflow orchestration frameworks focus on troubleshooting workflows after failures occur, not preventing failures proactively. Schema-only validation at phase boundaries misses semantic violations (keyword patterns) and compositional failures (cross-field logic, state-based checks), allowing infeasible hypotheses to propagate through expensive implementation phases before discovery. No existing research applies formal constraint validation methods to research pipeline phase transitions.

This work addresses the gap with contract-based validation: three-layer validation (schema + pattern + compositional contracts) forces explicit checking of constraints that schema-only validation cannot express. Each layer targets a distinct violation class: schemas validate structure (types, required fields), pattern validators check semantic content (forbidden keywords), and compositional contracts enforce cross-field logic (conditional dependencies) and state-based rules (membership in external registries). The layers are complementary---each catches violations the previous layer misses.

The contributions are:

\begin{enumerate}
\item First application of Design-by-Contract to semantic constraint validation in research workflows. Preconditions, postconditions, and invariants from code correctness (Eiffel, C++26 contracts) are adapted to constraint-preserving phase transitions in research pipelines.
\end{enumerate}

\begin{enumerate}
\item Three-layer validation architecture. A schema $\rightarrow$ pattern $\rightarrow$ contract hierarchy where each layer complements the previous, achieving 88\% violation detection at boundaries versus 40\% for schema-only (48 percentage point gap).
\end{enumerate}

\begin{enumerate}
\item Empirical validation on placeholder workflows. 100\% downstream failure reduction on a 100-case placeholder hypothesis corpus with four feasibility constraints on placeholder content with typed interfaces. External validity to substantive research is unproven and deferred to future work.
\end{enumerate}

\begin{enumerate}
\item Contract specification templates for rapid deployment. Reusable templates for semantic constraints (keyword blacklists), cross-field constraints (conditional membership), and state-based constraints (registry validation) enable adoption with minimal implementation cost (25-27 lines of code for the contract layer).
\end{enumerate}
